\DocumentMetadata{
  pdfversion=2.0,
  pdfstandard=ua-2,
  lang=en-US,
  tagging=on,
  tagging-setup={math/setup=mathml-SE,tabsorder=structure}
}
\documentclass[11pt,letterpaper]{article}
\usepackage[margin=0.68in,includefoot]{geometry}
\usepackage{newcomputermodern}
\usepackage{amsmath,amscd,mathtools}
\providecommand{\square}{\mdlgwhtsquare}
\usepackage[hidelinks]{hyperref}
\hypersetup{
  pdftitle={Logic PhD Exam, September 1992},
  pdfauthor={University of Florida Department of Mathematics},
  pdfsubject={Graduate mathematics examination},
  pdfdisplaydoctitle=true,
  bookmarksopen=true
}
\setcounter{secnumdepth}{0}
\setlength{\parindent}{0pt}
\setlength{\parskip}{0.28em}
\setlength{\emergencystretch}{3em}
\setlength{\leftmargini}{2.15em}
\setlength{\leftmarginii}{2.65em}
\setlength{\leftmarginiii}{3em}
\renewcommand{\labelenumi}{\textbf{\theenumi.}}
\pagestyle{empty}
\raggedbottom
\allowdisplaybreaks
\newenvironment{examproblems}[1]{%
  \begin{enumerate}%
  \setcounter{enumi}{#1}%
  \setlength{\topsep}{0.35em}%
  \setlength{\partopsep}{0pt}%
  \setlength{\itemsep}{0.48em}%
  \setlength{\parsep}{0.18em}%
}{\end{enumerate}}
\newenvironment{examsubparts}{%
  \begin{enumerate}%
  \setlength{\topsep}{0.18em}%
  \setlength{\partopsep}{0pt}%
  \setlength{\itemsep}{0.16em}%
  \setlength{\parsep}{0.1em}%
}{\end{enumerate}}
\newenvironment{examinstructions}{%
  \par\begingroup\itshape
}{%
  \par\endgroup\vspace{0.18em}
}
\makeatletter
\renewcommand\section{\@startsection{section}{1}{0pt}%
  {-1.25ex plus -.25ex minus -.1ex}%
  {0.55ex plus .1ex}%
  {\normalfont\large\bfseries}}
\renewcommand{\@maketitle}{%
  \newpage\null\vskip -1.2em%
  {\footnotesize\bfseries
    \href{https://www.ufl.edu/}{University of Florida}, \href{https://math.ufl.edu/}{Department of Mathematics}\par}%
  \vskip 0.18em%
  \tagstructbegin{tag=Title}%
    {\LARGE\bfseries\href{https://gma.math.ufl.edu/exams/logic-phd/}{Logic PhD Exam}, September 1992\par}%
  \tagstructend%
  \vskip 0.35em%
  \tagmcbegin{artifact}%
    \hrule height 1pt%
  \tagmcend%
  \vskip 0.55em%
}
\makeatother
\title{Logic PhD Exam, September 1992}
\author{}
\date{}
\begin{document}
\maketitle
\thispagestyle{empty}
\begin{examinstructions}
This examination is divided into 4 segments, with 3 questions in each segment. Answer any 2 questions from each segment.
\end{examinstructions}
\section{1. General Logic}
\begin{examproblems}{0}
\item
Prove that there is no single sentence~\(\phi\) in the language of ring theory such that a ring~\(R\) is a model of~\(\phi\) if and only if~\(R\) has characteristic 0.
\item
State and prove the compactness theorem for Predicate logic. You may assume the completeness theorem.
\item
State and sketch a proof of Gödel’s incompleteness theorem.
\end{examproblems}
\section{2. Model Theory}
\begin{examproblems}{0}
\item
Show that any two elementarily equivalent, countably saturated models are isomorphic.
\item
State and sketch a proof of the downward Löwenheim–Skolem Theorem.
\item
State the Omitting types theorem, and use it to show that any~\(\omega\)-complete theory of arithmetic has an~\(\omega\)-model.
\end{examproblems}
\section{3. Set Theory}
\begin{examproblems}{0}
\item
Show that a countable union of countable sets is countable. Explain any use of the axiom of choice.
\item
Show that \(2^\kappa = \kappa^\kappa\) for all infinite~\(\kappa\).
\item
Use Fodor’s theorem to prove that the diagonal intersection of a sequence of closed, unbounded sets is closed and unbounded.
\end{examproblems}
\section{4. Recursion Theory}
\begin{examproblems}{0}
\item
Show that a function \(\phi : \omega \to \omega\) is (partial) recursive iff its graph is recursively enumerable, and that if~\(\phi\) is total and recursive then its graph is recursive.
\item
Prove that there is a subset of~\(\omega\) which is recursively enumerable but not recursive.
\item
State the normal form theorem for partial recursive functions and use it to prove the enumeration theorem.
\end{examproblems}
\end{document}
